\documentclass[12pt,a4paper]{article}

\usepackage{cmap}
\usepackage[T2A]{fontenc}
\usepackage[utf8]{inputenc}
\usepackage[russian]{babel}
\usepackage{amsmath,amssymb}
\usepackage[a4paper,margin=2cm]{geometry}

\setlength{\parindent}{0pt}
\setlength{\parskip}{0.45em}
\sloppy

\begin{document}

\begin{center}
{\Large\bfseries Билет 13: вопросы для проверки знаний}
\end{center}

\noindent\fbox{%
  \begin{minipage}{\dimexpr\linewidth-2\fboxsep-2\fboxrule\relax}
  \normalfont
Степенные ряды с действительными членами. Почленное интегрирование
степенного ряда. Бесконечная дифференцируемость суммы степенного ряда на
интервале сходимости. Единственность представления функции степенным рядом.
Достаточные условия разложимости бесконечно дифференцируемой функции в
степенной ряд. Ряд Тейлора. Формула Тейлора с остаточным членом в
интегральной форме. Пример бесконечно дифференцируемой функции, не
разлагающейся в степенной ряд. Разложение в ряд Тейлора основных
элементарных функций:
\(e^x\), \(\cos x\), \(\sin x\), \(\ln(1+x)\), \((1+x)^\alpha\). Разложение в степенной
ряд функции \(e^z\) комплексного переменного \(z\).
  \end{minipage}%
}

\section*{Вопросы на понимание}

\begin{enumerate}
\item Почему радиус сходимости степенного ряда не дает полного ответа о сходимости в граничных точках интервала?
\item Что может пойти не так, если пытаться почленно дифференцировать степенной ряд прямо в граничной точке интервала сходимости?
\item Почему для почленного интегрирования достаточно равномерной сходимости на отрезке, лежащем строго внутри интервала сходимости?
\item Если два степенных ряда с одним и тем же центром задают одну функцию в окрестности центра, почему их коэффициенты уже не могут отличаться?
\item Чем утверждение формулы Тейлора с остаточным членом отличается от утверждения, что функция равна своему ряду Тейлора?
\item Почему бесконечная дифференцируемость функции сама по себе еще не гарантирует разложимость в степенной ряд?
\item Что именно показывает пример \(f(x)=e^{-1/x^2}\) при \(x\ne0\), \(f(0)=0\), о связи производных в точке и значений функции рядом с этой точкой?
\item Почему условие единой ограниченности всех производных в окрестности сильнее, чем просто существование всех производных?
\item Как интегральная форма остатка помогает понять, что остаток в формуле Тейлора стремится к нулю?
\item Почему разложение \(\ln(1+x)\) требует отдельного внимания к точке \(x=1\), хотя внутри \(|x|<1\) оно получается интегрированием геометрического ряда?
\item Что меняется в биномиальном ряде для \((1+x)^\alpha\), когда \(\alpha\) является неотрицательным целым числом?
\item Как из комплексного ряда для \(e^z\) естественно появляются вещественные разложения для \(\cos x\) и \(\sin x\)?
\end{enumerate}

\section*{Источники для сверки}

\texttt{target/answer\_question\_13.tex};
\texttt{IvGE\_1.pdf}, стр. 304--321.

\end{document}
